\let\negmedspace\undefined
\let\negthickspace\undefined
\documentclass[journal]{IEEEtran}
\usepackage[a5paper, margin=10mm, onecolumn]{geometry}
%\usepackage{lmodern} % Ensure lmodern is loaded for pdflatex
\usepackage{tfrupee} % Include tfrupee package

\setlength{\headheight}{1cm} % Set the height of the header box
\setlength{\headsep}{0mm}     % Set the distance between the header box and the top of the text

\usepackage{gvv-book}
\usepackage{gvv}
\usepackage{cite}
\usepackage{amsmath,amssymb,amsfonts,amsthm}
\usepackage{algorithmic}
\usepackage{graphicx}
\usepackage{textcomp}
\usepackage{xcolor}
\usepackage{txfonts}
\usepackage{listings}
\usepackage{enumitem}
\usepackage{mathtools}
\usepackage{gensymb}
\usepackage{comment}
\usepackage[breaklinks=true]{hyperref}
\usepackage{tkz-euclide} 
\usepackage{listings}
% \usepackage{gvv}                                        
\def\inputGnumericTable{}                                 
\usepackage[latin1]{inputenc}                                
\usepackage{color}                                            
\usepackage{array}                                            
\usepackage{longtable}                                       
\usepackage{calc}                                             
\usepackage{multirow}                                         
\usepackage{hhline}                                           
\usepackage{ifthen}                                           
\usepackage{lscape}
\begin{document}

\bibliographystyle{IEEEtran}
\vspace{3cm}

\title{7.4.40}
\author{EE25BTECH11033 - Kavin}
% \maketitle
% \newpage
% \bigskip
{\let\newpage\relax\maketitle}

\renewcommand{\thefigure}{\theenumi}
\renewcommand{\thetable}{\theenumi}
\setlength{\intextsep}{10pt} % Space between text and floats
\textbf{Question}:\\
Let a circle be given by $2x\brak{x-a}+y\brak{2y-b}=0$, $\brak{a\neq0, b\neq0}$. Find the condition on $a$ and $b$ if two chords,  each bisected by the $X$ axis, can be drawn to the circle from $\brak{a, \frac{b}{2}}$.\\    
\bigskip

\textbf{Solution}:\\
The equation of the circle is ,
\begin{align}
    2x\brak{x-a}+y\brak{2y-b}=0
\end{align}
\begin{align}
    \implies2x^2 - 2ax + 2y^2 - by = 0
\end{align}
\begin{align}
    \implies x^2 + y^2 - ax - \frac{b}{2}y = 0
\end{align}\\
\bigskip
The chords are drawn from the point $\vec{P}(a, \frac{b}{2})$ , which lies on the circle since\\
\begin{align}
a^2 + \brak{\frac{b}{2}}^2 - a(a) - \frac{b}{2}\brak{\frac{b}{2}} = a^2 + \frac{b^2}{4} - a^2 - \frac{b^2}{4} = 0
\end{align}


Let a chord start at $\vec{P}(a, \frac{b}{2})$ and end at another point on the circle, $\vec{Q}(x_q, y_q)$. The midpoint of this chord is $\vec{M} = \brak{\frac{a+x_q}{2}, \frac{b/2+y_q}{2}}.\\

Since this chord is bisected by the X-axis,
\begin{align}
\frac{b/2+y_q}{2} = 0 \implies \frac{b}{2} + y_q = 0 \implies y_q = -\frac{b}{2}
\end{align}
Since $\vec{Q}(x_q, -\frac{b}{2})$ must also lie on the circle,
\begin{align}
    x_q^2 + \brak{-\frac{b}{2}}^2 - a(x_q) - \frac{b}{2}\brak{-\frac{b}{2}} &= 0 \\
    x_q^2 + \frac{b^2}{4} - ax_q + \frac{b^2}{4} &= 0 \\
    x_q^2 - ax_q + \frac{b^2}{2} &= 0
\end{align}\\

The problem requires that two such chords can be drawn. This implies that the quadrateic equation for $x_q$ must have two distinct real roots.\\

Any monic quadratic equation of the form $\lambda^2 - S\lambda + P = 0$ can be interpreted as the \textbf{characteristic equation} of a $2 \times 2$ matrix $\vec{A}$, where $S$ is the trace and $P$ is the determinant of the matrix.
\begin{align}
\mydet{\vec{A} - \lambda \vec{I}} = \lambda^2 - \text{tr}(\vec{A})\lambda + \mydet{\vec{A}} = 0
\end{align}
The roots of this equation are the \textbf{eigenvalues} of the matrix $\vec{A}$.

Here,
\begin{itemize}
    \item  $\lambda = x_q$
    \item  $\text{tr}(\vec{A}) = a$
    \item  $\mydet{\vec{A}} = \frac{b^2}{2}$
\end{itemize}
A fundamental property of a $2 \times 2$ matrix is that it has \textbf{two distinct real eigenvalues} if and only if the discriminant of its characteristic equation is positive. This condition, in terms of its trace and determinant, is:
\begin{align}
(\text{tr}(\vec{A}))^2 - 4\mydet{\vec{A}} > 0
\end{align}
With $\text{tr}(\vec{A}) = a$ and $\mydet{\vec{A}} = \frac{b^2}{2}$, the condition becomes:
\begin{align}
    a^2 - 4\brak{\frac{b^2}{2}} &> 0
\end{align}
\begin{align}
    a^2 - 2b^2 &> 0 
\end{align}
\begin{align}
    a^2 &> 2b^2
\end{align}
This is the required condition on $a$ and $b$. \\
\bigskip

For $a = 2$ and $b = 1$,
\newpage
\begin{figure}[H]
\begin{center}
\includegraphics[width=0.6\columnwidth]{figs/fig.png}
\end{center}

\label{fig:Fig1}
\end{figure}
\end{document}


